\section{Actual module isomorphisms preserve the original model norm}

\begin{lemma}[Original module isomorphisms preserve genuine generic frames]
\label{lem:scheme-line-actual-generic-frame-naturality}
\lean{EnvelopeProperty.Berkovich.RelativeSchemeValuedPoint.centerFrameGenericIso_natural}
\leanok
\uses{lem:scheme-line-actual-center-frame}
For an actual original module isomorphism
$a:\mathcal M\simeq\mathcal N$ and an actual valuative frame $e$ of
$\mathcal N$, composing its actual center pullback with $e$ gives a
frame of $\mathcal M$. Its actual generic frame is the original
generic pullback of $a$ followed by the actual generic frame of $e$.
\end{lemma}
\begin{proof}
The actual generic-comparison isomorphism is the native pullback
composition followed by the native equality comparison from the
actual valuative triangle. It is consequently natural in every
original module morphism. Its inverse naturality on $a$ moves the
actual generic pullback of $a$ through that inverse comparison.
Functoriality of the original center and generic pullbacks then
identifies the two displayed composites of actual frames.
\end{proof}

\begin{lemma}[An actual whole original line isomorphism is an isometry]
\label{lem:scheme-line-actual-model-isomorphism-isometry}
\lean{EnvelopeProperty.Berkovich.RelativeSchemeValuedPoint.modelLineNorm_mapIso}
\leanok
\uses{lem:scheme-line-actual-generic-frame-naturality,def:scheme-line-actual-model-norm}
Every actual original module-sheaf isomorphism between invertible
sheaves is an isometry of their genuine proper-model norms on the
same original residue-field pullback line fibers.
\end{lemma}
\begin{proof}
Choose a derived actual valuative frame $e$ of the target. Compute
its original model norm in that frame. For the source use the actual
center pullback of the given isomorphism followed by $e$; frame
independence permits this frame. The preceding genuine generic-frame
naturality identifies its coefficient with the target coefficient of
the actual original generic image of the section. Their original
residue-field valuations, hence the two model norms, are equal.
\end{proof}

\begin{lemma}[The actual structure-module model norm is the original scalar norm]
\label{lem:scheme-line-actual-model-unit-norm}
\lean{EnvelopeProperty.Berkovich.RelativeSchemeValuedPoint.modelLineNorm_unit}
\leanok
\uses{def:scheme-line-actual-model-norm}
For the actual structure-module line, the genuine proper-model norm
is the original residue-field valuation after the canonical original
structure-module pullback comparison.
\end{lemma}
\begin{proof}
Use the actual structure-unit pullback comparison at the genuine
valuative center as the actual center frame. Monoidality of the
original structure-unit comparison for successive pullbacks identifies
its generic frame with the canonical structure-unit pullback
comparison at the original residue point: the native comparison
satisfies the true composition equation, and the native equality
comparison satisfies the true unit equation along the valuative
triangle. Cancel these actual comparison isomorphisms. By the
definition of the original frame norm the result is exactly the
original residue-field scalar valuation, as stated.
\end{proof}
